%!TeX root=MemoriaTFG.tex
%
% ============================================================================
% NOTA PARA EL AGENTE QUE INTEGRA EL CAPITULO
% ----------------------------------------------------------------------------
% 1) Este capitulo es el puente investigacion -> producto: documenta el
%    PROTOTIPO DESPLEGADO (dermapixel.eu), distinto del capitulo de
%    EVALUACION (\ref{cap:dermapixel}) y del Anexo arquitectonico M1--M11
%    (\ref{cap:anexoh}). Reusa esas referencias en vez de duplicar contenido.
% 2) Verificar que existen las labels referenciadas antes de compilar:
%    cap:resultados, cap:dermapixel, cap:anexoh, cap:metodos, sec:zs,
%    sec:dermapixel-m11, sec:spanderm-clip, sec:linea-prospectivo,
%    sec:alcance. Si alguna cambia de nombre, ajustar \ref aqui.
% 3) FIGURAS: los \includegraphics estan COMENTADOS para que compile sin las
%    imagenes. Capturar los pantallazos indicados en cada bloque
%    "% CAPTURAR:", guardarlos en figures/ con el nombre indicado y
%    descomentar la linea. Lista completa al final de la introduccion.
% 4) Cifras y nombres de modulo verificados contra codigo real
%    (pipeline.py, spark_adapter.py, docker-compose.yml, models.py,
%    ontology.csv). No introducir cifras sin fuente.
% ============================================================================

\chapter{El prototipo DermApIxel: del banco de pruebas al asistente clínico}
\label{cap:prototipo}

Los capítulos anteriores responden a una pregunta de investigación: qué
modelos fundacionales sirven, y para qué tarea, sobre material
dermatológico ---incluido el material clínico en castellano de
DermapixelAI~1.0 (capítulo~\ref{cap:dermapixel})---. Este capítulo responde
a otra complementaria: \emph{cuánto de lo estudiado sobrevive al pasar de
un \texttt{notebook} de evaluación a un sistema que un dermatólogo puede
usar de verdad}. No todo lo que mejora una métrica merece desplegarse, y el
prototipo es donde esa distinción se vuelve operativa. No es objeto
experimental del trabajo (sección~\ref{sec:alcance}) ---la defensa empírica
descansa en los modelos del capítulo~\ref{cap:resultados} bajo protocolos
estándar---: lo que aporta es la \emph{traslación}, materializar esos
modelos en una interfaz clínica multicanal y tomar las decisiones de
seguridad clínica que un benchmark no obliga a tomar pero un sistema en
producción sí.

La arquitectura interna de los once módulos (M1--M11) se detalla en el
repositorio reproducible (\ref{cap:repo}, sección~\ref{sec:repo-prototipo});
aquí se describe su papel dentro del producto y las capas que el
repositorio no cubre: el asistente conversacional, los tres canales (web,
WhatsApp y voz) y las herramientas externas. A fecha de cierre, el sistema
está en producción \emph{end-to-end} en un servidor público
(\texttt{dermapixel.eu}), con la Dra.~R.~Taberner como evaluadora clínica de
referencia; sigue sin validarse sobre cohorte prospectiva ni integrarse en
infraestructura hospitalaria (sección~\ref{sec:proto-estado}), su línea de
continuación inmediata (sección~\ref{sec:linea-prospectivo}).

% ---------------------------------------------------------------------------
% PLAN DE FIGURAS (pantallazos a capturar). Capturar en dermapixel.eu con
% una cuenta de rol experto y un caso de ejemplo no sensible:
%   figures/proto_consenso.png   -> Vista de analisis: banner de consenso
%                                   multi-modulo + veredicto + urgencia.
%   figures/proto_7puntos.png    -> Tab 7-POI sobre imagen DERMATOSCOPICA
%                                   (panel visible) con los 7 criterios.
%   figures/proto_seg.png        -> Tab SEG: overlay de mascara sobre lesion.
%   figures/proto_zeroshot.png   -> Tab Z-SHOT: comparacion abierta con tags
%                                   L1/L2/L3 y badge de cascada principal.
%   figures/proto_ragES.png      -> Tab RAG-I (M4-bis): casos similares del
%                                   archivo de la Dra. Taberner + voto.
%   figures/proto_chat_web.png   -> Asistente conversacional web con citas
%                                   al blog Dermapixel.
%   figures/proto_whatsapp.png   -> Conversacion WhatsApp (movil) con analisis
%                                   de imagen.
%   figures/proto_voz.png        -> Widget de voz activo + captura cross-device
%                                   (QR a la PWA movil).
%   figures/proto_historial.png  -> Historial paginado (grid de thumbnails).
%   figures/proto_arquitectura.png-> (Opcional) diagrama de microservicios.
% ---------------------------------------------------------------------------

\section{De la evaluación al producto: qué se traslada}
\label{sec:proto-mapeo}

La tabla~\ref{tab:proto-mapeo} resume la correspondencia entre cada bloque
estudiado en el trabajo y su materialización en el prototipo. La columna
\emph{Estado} distingue tres situaciones: desplegado y expuesto al
clínico, desplegado como capa interna (no visible directamente), y
estudiado pero \emph{no} desplegado por una decisión justificada.

\begin{table}[htbp]
\centering
\footnotesize
\setlength{\tabcolsep}{4pt}
\renewcommand{\arraystretch}{1.05}
\caption{Trazabilidad estudio~$\to$~prototipo. Los nombres de módulo
(M1--M11) se conservan tal como aparecen en \texttt{pipeline.py} para
preservar la correspondencia código--documento (repositorio
reproducible, sección~\ref{sec:repo-prototipo}).}
\label{tab:proto-mapeo}
\adjustbox{max width=\textwidth}{%
\begin{tabular}{lll}
\hline
Estudiado en el trabajo & Materialización en el prototipo & Estado \\
\hline
Clasificación supervisada PanDerm (\ref{cap:resultados}) & M1 (HAM10000) / M7 (jerárquico \emph{merged43}+TTA) & Desplegado, visible \\
Dermapixel R0, head L2 castellana (\ref{cap:dermapixel}) & M9 (PanDerm + LoRA $r{=}16$) & Desplegado, visible \\
Segmentación promptable (\ref{cap:resultados}) & M2/segmentación + corrección experta & Desplegado, visible \\
Interpretabilidad por SAE + SkinCon (\ref{cap:resultados}) & M3 (32 conceptos por imagen) & Desplegado, visible \\
\emph{Seven-Point Checklist} multitarea (\ref{cap:resultados}) & M10 + puntuación Argenziano (capa \emph{backend}) & Desplegado, visible \\
Zero-shot multimodal DermLIP (\ref{sec:zs}) & M6 (clase abierta) + glosario ES$\to$EN & Desplegado, visible \\
Recuperación visual densa (\ref{cap:resultados}) & M4 (Derm1M) / M4-bis (DermapixelAI~ES) & Desplegado, visible \\
Comparación con LLM multimodales (\ref{cap:resultados}) & Razonamiento clínico estructurado (M5) & Desplegado, visible \\
Ensemble por complementariedad (\ref{sec:repo-prototipo}) & M11 + banner de consenso de 5 módulos & Desplegado, visible \\
Ontología jerárquica L1/L2/L3 (\ref{cap:metodos}) & Coherencia de niveles + traducción EN$\to$ES & Desplegado, interno \\
Texto clínico como contexto (\ref{cap:dermapixel}) & RAG conversacional sobre el archivo Dermapixel & Desplegado, visible \\
Dermapixel R0, rama contrastiva (\ref{sec:spanderm-clip}) & SpanDerm-CLIP (búsqueda textual ES) & \textbf{No desplegado} \\
\hline
\end{tabular}}
\end{table}

La lectura transversal de la tabla expone el mensaje central del
capítulo: el prototipo no es una demostración de un único modelo, sino la
integración de \emph{varias} líneas del trabajo, cada una en el papel para
el que el banco de pruebas la encontró útil. El clasificador jerárquico
M7 aporta cobertura sobre las clases del archivo de la Dra.~R.~Taberner; la
adaptación ligera M9 aporta el vocabulario castellano; el zero-shot M6
aporta apertura a clases no vistas; la recuperación M4-bis aporta casos
comparables del propio archivo; y el razonamiento M5 los integra en una
explicación textual. Ninguno sustituye al juicio clínico: el prototipo se
diseña como apoyo, no como decisor.

\section{Arquitectura del sistema desplegado}
\label{sec:proto-arq}

El prototipo evolucionó desde el despliegue local documentado en el
repositorio reproducible (\ref{cap:repo}, sección~\ref{sec:repo-prototipo};
PC Windows~$+$~DGX~Spark sobre red LAN, estado 2026-04) hacia un despliegue
en servidor accesible públicamente. El cambio
no es solo de ubicación: introduce la necesidad de autenticación,
cumplimiento de protección de datos y los tres canales de interacción
descritos en la sección~\ref{sec:proto-asistente}.

\subsection{Topología de microservicios}
\label{sec:proto-micro}

El \emph{backend} se organiza como microservicios en contenedores
(Docker Compose) repartidos en dos redes: una expuesta al exterior (proxy
inverso y salida a APIs externas) y otra interna sin acceso directo a
Internet, por mínima exposición. La tabla~\ref{tab:proto-services} resume
los servicios principales.

\begin{table}[htbp]
\centering
\footnotesize
\setlength{\tabcolsep}{4pt}
\renewcommand{\arraystretch}{0.98}
\caption{Servicios del \emph{backend} de producción. Los servicios MCP
(\emph{Model Context Protocol}) exponen herramientas externas al
orquestador conversacional (sección~\ref{sec:proto-mcp}).}
\label{tab:proto-services}
\begin{tabular}{lll}
\hline
Servicio (contenedor) & Tecnología & Función \\
\hline
\texttt{auth-svc} (\texttt{dpai-auth})    & FastAPI       & Identidad, 2FA, JWT RS256, auditoría \\
\texttt{blog-svc} (\texttt{dpai-blog})    & FastAPI       & API principal, RAG, asistente, análisis \\
\texttt{whatsapp-svc} (\texttt{dpai-wha.})& FastAPI       & \emph{Webhook} Meta Cloud API \\
\texttt{spark-svc} (\texttt{dpai-spark})  & FastAPI       & Proxy resiliente a la inferencia GPU \\
\texttt{nginx} (\texttt{dpai-nginx})      & Nginx         & Proxy inverso, TLS, ficheros estáticos \\
MySQL (\texttt{dpai-mysql})               & MySQL $8.0$   & Persistencia (BD \texttt{dpai\_auth}, \texttt{dpai\_blog}) \\
Redis (\texttt{dpai-redis})               & Redis $7$     & Caché y estado conversacional efímero \\
RabbitMQ (\texttt{dpai-rabbitmq})         & RabbitMQ $3.13$& Mensajería asíncrona \\
Qdrant (\texttt{dpai-qdrant})             & Qdrant        & Índice vectorial del RAG conversacional \\
MCP \texttt{pubmed/drugs/ontology/spark}  & FastMCP       & Herramientas externas (sección~\ref{sec:proto-mcp}) \\
\hline
\end{tabular}
\end{table}

Las dos bases de datos están separadas por servicio (identidad y datos
clínicos), sin enlace directo entre ellas: la relación usuario--estudio se
comprueba en la aplicación. El aislamiento cuesta algún paso extra en
operaciones como el borrado, a cambio de limitar el alcance de un posible
incidente sobre datos de salud.

\subsection{Separación cliente / inferencia GPU}
\label{sec:proto-spark}

La inferencia pesada (modelos M1--M11) corre en el servidor GPU propio
(DGX~Spark), separado del \emph{backend} y accesible solo por una red
privada (Tailscale). El \emph{backend} no llama a la GPU directamente,
sino a través del servicio \texttt{spark-svc}, que la convierte en una
dependencia tolerante a fallos: reintenta los errores, deja de insistir si
la GPU falla repetidamente y guarda en caché lo ya calculado. Así, una
caída momentánea de la GPU degrada el servicio con elegancia ---el estudio
queda \emph{en proceso} y un proceso de limpieza lo marca como fallido
pasado un tiempo--- en lugar de propagar el error al clínico.

El acceso a la inferencia se expone además como herramienta del asistente
conversacional mediante un servicio MCP dedicado
(\texttt{mcp-spark}), que publica las operaciones
\texttt{analyze\_dermatology\_image}, \texttt{zero\_shot\_classify},
\texttt{find\_similar\_cases} y \texttt{extract\_clinical\_concepts}. De
este modo el mismo \emph{pipeline} de diagnóstico es accesible tanto desde
la vista de análisis (uso directo) como desde una conversación en
lenguaje natural (uso orquestado por el LLM).

\subsection{Persistencia, trazabilidad y rendimiento}
\label{sec:proto-persistencia}

Cada estudio diagnóstico se persiste en la tabla \texttt{derm\_studies}
con su imagen (almacenamiento de objetos cifrado en reposo), un
\emph{thumbnail} pre-renderizado, el canal de origen
(\texttt{web}/\texttt{whatsapp}/\texttt{voice}), la salida cruda de la
inferencia, el razonamiento generado y el tipo de imagen detectado
(clínica o dermatoscópica). Las validaciones del experto, el
\emph{feedback} y las sesiones de anotación se enlazan a este registro y
alimentan el aprendizaje activo (sección~\ref{sec:proto-anotacion}). Toda
acción sensible se registra en una tabla de auditoría con la IP
\emph{hasheada} (SHA-256, sin IP en claro), conforme a protección de
datos; \texttt{auth-svc} mantiene además una bitácora de auditoría
encadenada criptográficamente (\emph{hash} encadenado) orientada al
Esquema Nacional de Seguridad.

Una optimización con impacto directo en la experiencia fue generar las
miniaturas del historial al subir la imagen y servirlas por una URL
firmada y temporal, en lugar de incrustar cada imagen dentro de la
respuesta: el historial pasó a cargar mucho más rápido y dejó de fallar al
pedir muchas miniaturas a la vez.

\section{El pipeline diagnóstico integrado}
\label{sec:proto-pipeline}

La vista de análisis expone el resultado de los módulos M1--M11 en
pestañas y, sobre ellas, un \emph{banner} de consenso que resume la
lectura de malignidad de cinco clasificadores independientes (M1, M7, M9,
sondeo lineal SigLIP y M10) más un sexto indicador con la ontología del
vecino \emph{top-1} castellano (M4-bis). Cuando dos o más módulos votan
melanoma con probabilidad alta, el sistema emite un aviso explícito de
derivación urgente. El diseño replica el patrón clínico de segunda
opinión: no se confía en un único clasificador, se busca convergencia
entre criterios que explotan representaciones parcialmente distintas
---justamente la complementariedad demostrada en la
sección~\ref{sec:repo-prototipo}---.

Las principales vistas de la interfaz clínica del prototipo se recogen, en
miniatura, en las figuras~\ref{fig:proto-interfaz}
y~\ref{fig:proto-interfaz2} (al final del capítulo).

El veredicto que se muestra al clínico no es la salida de un único modelo
sino el ensemble ponderado M11, cuya jerarquía L1/L2/L3 se hace coherente
\emph{top-down} (si L1 indica patología tumoral y L3 indica melanoma, L2
no puede indicar una categoría benigna) y cuyas etiquetas se unifican
mediante un \emph{mapping} EN$\to$ES que colapsa las duplicaciones entre
el M7 (que emite la ontología en inglés) y los módulos castellanos M9 y
M4-bis. La coherencia de niveles se apoya en la ontología jerárquica
validada con la Dra.~R.~Taberner (4~categorías L1, 43~subcategorías L2 y
$\sim$367 diagnósticos L3; repositorio reproducible,
sección~\ref{sec:repo-prototipo}).

\subsection{Tres modos de interpretación, una misma red de seguridad}
\label{sec:proto-modos}

Una decisión de diseño central es que el clínico elige \emph{cuánto}
quiere que la IA interprete el caso, mediante un selector con tres modos
que comparten el mismo análisis de base pero dan distinto peso al
razonamiento clínico estructurado (M5):

\begin{itemize}
  \item \textbf{Determinista.} El veredicto procede únicamente de los
  clasificadores (el ensemble M11); el módulo de razonamiento no
  interviene. Es reproducible ---la misma entrada da siempre el mismo
  resultado--- y está siempre disponible, también como respaldo si el
  razonamiento falla.
  \item \textbf{Híbrido (IA-asistido)}, el modo por defecto. Parte del
  veredicto determinista y lo \emph{refina} con las señales de M5: solo
  cambia el diagnóstico principal si el razonamiento propone otro
  \emph{con alta confianza}, y puede aportar el resumen narrativo y el
  diagnóstico diferencial ordenado.
  \item \textbf{Interpretación IA máxima.} El razonamiento de M5 lleva la
  voz cantante: el diagnóstico principal se sustituye por el que propone y
  el \emph{pipeline} de modelos queda como contexto. Si M5 no ofrece un
  diagnóstico claro, el sistema vuelve al determinista.
\end{itemize}

Sobre los tres modos rige una misma red de seguridad innegociable:
\textbf{la urgencia nunca baja por acción de la IA} ---el razonamiento
puede \emph{subir} el nivel de alarma ante una lesión sospechosa, nunca
rebajarlo--- y el veredicto determinista se conserva siempre como base
auditable. Así el clínico gana interpretabilidad sin que la capa
generativa pueda relajar una señal de seguridad, una decisión que un banco
de pruebas no obliga a tomar pero un sistema en producción sí.

% CAPTURAR: el selector de modo (Determinista / Híbrido / Interpretación
% IA máxima) puede mostrarse en la misma figura de la vista de análisis
% (fig:proto-consenso) o como recorte propio.

\subsection{La Lista de los 7 puntos: modelo de concepto y capa clínica}
\label{sec:proto-7pc}

El módulo M10 detecta por imagen los siete criterios dermatoscópicos del
\emph{Seven-Point Checklist} más una \emph{head} binaria de melanoma. La
separación de responsabilidades es deliberada: M10 solo emite las
probabilidades de cada criterio (una \emph{softmax} por concepto), y la
puntuación ponderada de Argenziano ---los tres criterios mayores valen
$2$, los menores $1$, melanoma a partir de $3$--- la calcula la capa
clínica del \emph{backend}, no la red, de modo que el esquema de puntuación
se puede auditar y ajustar sin reentrenar el modelo. La versión inicial
sumaba por conteo de criterios; el paso a puntuación ponderada es
conservador en seguridad: todos los pesos son~$\geq 1$, así que nunca se
pierde una alarma y se gana sensibilidad en las lesiones cargadas de
criterios mayores.

Una segunda decisión de seguridad es \emph{cuándo} mostrar el panel. La
Lista de los 7 puntos es dermatoscópica; aplicarla a una fotografía clínica
daría criterios falsos. Por eso solo se muestra sobre imagen
dermatoscópica, con doble comprobación ---el detector de modalidad (M0) y,
si hace falta, el experto--- y con un \emph{gate} \emph{fail-closed}: ante
duda, el panel no se aplica. La corrección del experto se guarda como
verdad de referencia para reentrenar (aprendizaje activo).

\subsection{Segmentación con corrección experta}
\label{sec:proto-seg}

La segmentación de la lesión, estudiada en el capítulo~\ref{cap:resultados}
como tarea promptable, se traslada al prototipo como una máscara
automática que el clínico ve superpuesta sobre la imagen y, si procede,
corrige. La capa de anotación registra explícitamente el \emph{origen} de
cada máscara aprobada ---máscara aceptada automáticamente, corregida a
partir de una propuesta automática, o dibujada desde cero--- lo que
permite distinguir, de cara al reentrenamiento, qué casos requirieron
intervención humana y cuáles no.

\subsection{Zero-shot abierto y el idioma del espacio latente}
\label{sec:proto-zeroshot}

El módulo M6 (DermLIP) permite comparar una imagen contra una lista
arbitraria de descripciones, complementando a los clasificadores cerrados
con apertura a clases no previstas. Su integración en el prototipo
documenta un hallazgo operativo con valor metodológico: \emph{la
sensibilidad de un modelo contrastivo al idioma de la consulta}. DermLIP
se entrenó con texto biomédico en inglés; al recibir nombres de
diagnóstico en castellano sin traducir, el \emph{ranking} se degradaba
hasta producir ordenaciones clínicamente erróneas (un carcinoma basocelular
quedaba por debajo de un melanoma en casos donde no debía). La causa no era
el modelo sino el cableado: faltaba la traducción ES$\to$EN para los
términos comunes. La corrección ---un glosario clínico ampliado mediante un
generador reproducible y marcado para revisión de la Dra.~R.~Taberner antes
de su validación--- restauró el comportamiento esperado.

Este episodio refuerza la disciplina de validación que recorre todo el
trabajo: una métrica de cabecera puede esconder un artefacto de régimen, y
conviene reproducirla en las condiciones reales de uso antes de construir
sobre ella (sección~\ref{sec:proto-noproducto}).

\section{El asistente conversacional multicanal}
\label{sec:proto-asistente}

La aportación que más distingue al prototipo de un clasificador de
imágenes es el asistente conversacional anclado al archivo de la
Dra.~R.~Taberner. Aquí se materializa la segunda vía que el trabajo
defiende como poco explorada (sección~\ref{sec:proto-mapeo}): el texto
clínico como contexto.

\subsection{RAG sobre el archivo Dermapixel}
\label{sec:proto-rag}

El asistente responde con \emph{retrieval-augmented generation} (RAG)
sobre el archivo del blog Dermapixel: busca los textos más relevantes
---combinando coincidencia de palabras y similitud semántica--- y, antes
de responder, comprueba si lo recuperado basta para hacerlo, para
responder con cautela o para declinar. El modelo por defecto es
\texttt{gpt-4o-mini} (con Claude Sonnet conmutable para comparar), y cada
respuesta incluye citas a las entradas del blog que la sustentan, para que
el clínico verifique la fuente.

\subsection{Tres canales: web, WhatsApp y voz}
\label{sec:proto-canales}

El asistente se expone por tres canales con una garantía explícita de
\emph{paridad}: una misma imagen y un mismo contexto producen el mismo
veredicto determinista con independencia del canal.

\begin{itemize}
  \item \textbf{Web.} Conversación en la propia aplicación, con acceso a la
  vista de análisis, al historial y a la validación experta.
  \item \textbf{WhatsApp.} Integración con la API Cloud de Meta mediante
  \emph{webhook} verificado por firma. Soporta conversación multi-turno
  (estado efímero en Redis) y análisis de imagen clínica, con vinculación
  del número a una cuenta autenticada y un límite diario de análisis por
  usuario. El número se almacena \emph{hasheado}, nunca en claro.
  \item \textbf{Voz.} Conversación hablada en tiempo real (API Realtime de
  OpenAI); el \emph{backend} solo entrega una credencial temporal de
  sesión. Durante la conversación el asistente puede pedir una foto: el
  usuario la arrastra o, desde otro dispositivo, escanea un código QR que
  abre una página de captura con la cámara; la imagen entra en el mismo
  \emph{pipeline} de análisis.
\end{itemize}

El canal de voz destapó problemas que un banco de pruebas no revela
---por ejemplo, un bucle de eco en el altavoz del iPhone, en el que el
asistente captaba su propia voz y se reiniciaba---, resueltos caso por
caso. Son la clase de problema que solo aparece al desplegar.

\subsection{Herramientas externas (MCP)}
\label{sec:proto-mcp}

El asistente no se limita al blog: dispone de herramientas externas (vía
\emph{Model Context Protocol}) que el LLM usa cuando hacen falta
---consultar PubMed, la información de medicamentos de la AEMPS y la
ontología clínica---. El acceso sigue una lista blanca: solo se permiten
las herramientas explícitamente autorizadas.

\subsection{Guardarraíles clínicos}
\label{sec:proto-guardrails}

El asistente incorpora una capa de seguridad con varios mecanismos:
validación del \emph{prompt} de entrada contra patrones de inyección y de
extracción de instrucciones; validación de la respuesta para evitar fugas
de secretos o de detalles de infraestructura; reglas clínicas explícitas
acordadas con la Dra.~R.~Taberner (preferencia por el principio activo
frente al nombre comercial salvo petición explícita; distinción cuidadosa
entre entidades con nombres parcialmente solapados); y un disparador de
urgencia (\emph{``llame al 112''}) que nunca es suprimido por los
guardarraíles de alcance ---la seguridad prima sobre la restricción de
tema---. Un principio transversal verificado en el trabajo es que una
inyección en el contexto clínico \emph{no puede} elevar la urgencia del
veredicto determinista: el texto libre del usuario no altera la salida de
los clasificadores.

\section{Validación experta y aprendizaje activo}
\label{sec:proto-anotacion}

El prototipo cierra el ciclo con el experto: el dermatólogo puede validar
cada análisis (correcto/parcial/incorrecto), corregir el diagnóstico por
niveles, corregir la modalidad de la imagen y, en el estudio de anotación,
corregir la máscara de segmentación y registrar los conceptos presentes.
Estas correcciones se persisten como verdad de referencia y constituyen la
materia prima de la validación clínica prospectiva planificada
(sección~\ref{sec:linea-prospectivo}).

\section{Ramas estudiadas no desplegadas}
\label{sec:proto-noproducto}

La honestidad metodológica exige documentar también lo que \emph{no} se
desplegó. La rama contrastiva de Dermapixel R0 (SpanDerm-CLIP,
sección~\ref{sec:spanderm-clip}), concebida como buscador textual en
castellano (texto~$\to$~imagen), \emph{no} se promovió a producción. Su
métrica de cabecera más favorable corresponde a un régimen concreto
(partición de validación, descripciones clínicas largas, un \emph{epoch}
específico) que no generaliza: en el régimen de uso real ---consulta corta
del usuario, conjunto de test no visto, archivo completo--- la recuperación
cae a un nivel cercano al azar. La partición es limpia (verificada
disjunta por caso e imagen), lo que descarta la contaminación como causa.
La conclusión honesta es que queda \emph{demostrada la viabilidad técnica
del espacio dual texto-imagen} (la arquitectura entrena, converge y los
brazos comparten un espacio recuperable bajo descripciones ricas), pero
\emph{no la capacidad de recuperación útil} en régimen real. La rama
permanece no desplegada; la búsqueda visual en producción la cubre M4-bis
(imagen~$\to$~imagen).

Este caso, junto con el del idioma del zero-shot
(sección~\ref{sec:proto-zeroshot}), motiva una práctica de aseguramiento
de calidad que el trabajo eleva a método: \emph{revalidar toda métrica de
cabecera bajo condiciones realistas de despliegue} (datos no vistos,
consulta corta, conjunto completo) y reportar con franqueza cuando no
generaliza. Dos resultados negativos se convierten así en una garantía de
rigor.

\section{Estado de despliegue y limitaciones}
\label{sec:proto-estado}

A fecha de cierre, el sistema está en producción \emph{end-to-end} sobre
\texttt{dermapixel.eu}, con los módulos M1--M11 operativos, el asistente
conversacional en sus tres canales y la Dra.~R.~Taberner como evaluadora
clínica de referencia. Este salto a producción no es solo técnico: el
proyecto fue distinguido con la \emph{Beca de Innovación e Inteligencia
Artificial} de la Academia Española de Dermatología y Venereología (AEDV),
dotada con $10\,000$~euros \emph{para su puesta en producción}
(53.\textsuperscript{er} Congreso Nacional de Dermatología y Venereología,
Maspalomas, mayo de 2026), lo que financia precisamente la consolidación
del prototipo como servicio clínico y es, en sí mismo, un aval externo de
su utilidad. No está validado sobre cohorte prospectiva real ni
integrado en la infraestructura del centro hospitalario. Las limitaciones y
condiciones para el despliegue clínico pleno ---validación prospectiva,
integración con sistemas de información hospitalaria, ponderación del
consenso calibrada con datos reales, y certificación como producto
sanitario--- se detallan como líneas de continuación inmediatas en la
sección~\ref{sec:linea-prospectivo}. El prototipo, en su estado actual,
demuestra la viabilidad operativa de integrar los hallazgos del trabajo en
un sistema accesible y usable, y constituye el entorno previsto para esa
validación.

\begin{figure}[H]
\centering
\newcommand{\protothumb}[2]{%
  \begin{minipage}[t]{0.48\textwidth}\centering
  \adjustbox{max width=\linewidth,max height=2.9cm,frame=0.4pt}{\includegraphics{#1}}\\[1pt]
  {\scriptsize #2}
  \end{minipage}}
\protothumb{proto_consenso.png}{(a) Veredicto y \emph{banner} de consenso}\hfill
\protothumb{proto_ensemble.png}{(b) Consenso del ensemble (M11)}

\vspace{2.5mm}
\protothumb{proto_seg.png}{(c) Segmentación: MedSAM2 $+$ U-Net}\hfill
\protothumb{proto_ragES.png}{(d) Casos similares del archivo (M4-bis)}

\vspace{2.5mm}
\protothumb{proto_7puntos.png}{(e) Lista de los 7 puntos (M10)}\hfill
\protothumb{proto_zeroshot.png}{(f) Zero-shot abierto (M6, DermLIP)}
\caption{Vista de análisis de DermApIxel ---el análisis de una imagen---:
(a)~el veredicto con \emph{banner} de consenso (aquí un caso
\emph{discordante} que dispara aviso de revisión); (b)~el consenso del
ensemble (M11); (c)~la segmentación con MedSAM2 y U-Net; (d)~los casos
similares del archivo de la Dra.~R.~Taberner (M4-bis, imagen$\to$imagen);
(e)~la Lista de los 7 puntos (M10, solo sobre dermatoscopia); y (f)~el
zero-shot abierto (M6, DermLIP).}
\label{fig:proto-interfaz}
\end{figure}

\begin{figure}[H]
\centering
\newcommand{\protothumb}[2]{%
  \begin{minipage}[t]{0.48\textwidth}\centering
  \adjustbox{max width=\linewidth,max height=3.2cm,frame=0.4pt}{\includegraphics{#1}}\\[1pt]
  {\scriptsize #2}
  \end{minipage}}
\protothumb{proto_razonamiento.png}{(g) Razonamiento clínico (M5)}\hfill
\protothumb{proto_chatweb.png}{(h) Asistente web (citas al blog)}

\vspace{2.5mm}
\protothumb{proto_whatsapp.png}{(i) WhatsApp (análisis en móvil)}\hfill
\protothumb{proto_annotation.png}{(j) Annotation Studio (conceptos)}

\vspace{2.5mm}
\protothumb{proto_historial.png}{(k) Historial de casos}\hfill
\protothumb{proto_datasets.png}{(l) Datasets de anotación}
\caption{Asistente conversacional y ciclo de datos de DermApIxel:
(g)~el razonamiento clínico educativo (M5); (h, i)~el asistente
conversacional anclado al blog, accesible por web y WhatsApp; y
(j)--(l)~el ciclo de anotación experta: \emph{Annotation Studio} de
conceptos, historial de casos y gestión de \emph{datasets} de anotación.}
\label{fig:proto-interfaz2}
\end{figure}
